\section{Introduction}
When a model writes code, reports a result or prepares a dataset, someone has to check it, and
increasingly that someone is another model. CrossAudit, the system every measurement here passes
through (\S\ref{sec:system}), is built on two rules for doing that. The model that audits must come
from a different vendor than the model that wrote the work, and what a deterministic check can verify
is verified by code, not by a model's opinion. An auditor of this kind reads work it did not write
and says whether to block it. We ask how much it catches, what decides the limit, and whether the
arrangement helps where checking matters most, in science.

A reviewer's score on a benchmark mixes at least four things: whether the model can see a defect,
whether it judges the defect worth blocking on, whether the benchmark's hidden tests define the
behaviour it would have to see, and how many times it is allowed to look. Work on critic models
improves the first two and reports the aggregate gain. We vary these factors one at a time where the
design allows it, against a truth no model supplies: a hidden test suite executed by an interpreter.
Some contrasts cannot be isolated: changing the reading model also changed its sampling, and the
role of the specification is a post hoc association.

\textbf{The limit on general-purpose code} (\S\ref{sec:ceiling}). Eight independent readings of the
shipped cross-vendor auditor reach 30.0\% union recall [20.0, 40.7] on defective code at 16.0\%
[10.1, 22.3] on correct code, and the eighth reading still adds 1.9 points, so returns diminish
without a plateau. A stronger same-vendor model, sampled differently, moves recall by $-26.4$
points [$-37.3$, $-15.6$] while flagging 3.3\% of correct code against 16.0\%, and the contrast
reverses sign under a weaker severity rule. One added rulebook sentence moves flags by $+26.8$ points
[13.6, 40.4] on an exploratory arm, at $+12.5$ points on correct code (an interval,
[$-3.6$, $+26.8$], that does not exclude zero), without measurably moving outcomes. Of 57 defects no
reading flagged, two model raters judge 44 to be behaviour the specification never determined (post
hoc). A recall figure is therefore contingent on auditor route, rulebook, decision rule and benchmark,
and is not portable without its false-positive rate.

\textbf{Auditing science} (\S\ref{sec:ai4s}). We then take the system into science, with one
preregistered study for each job a scientist would hand an auditor. On reported results, the
deterministic check that every number traces to its source flagged all 49 inconsistencies among a
report, its results file and its run log, and none of 25 results that were wrong consistently in all
three; the model auditor flagged 23 of those 25, 21 with a sound finding about the fault (post hoc).
On data, a validator written from the dataset's documentation and the auditor flagged different
faulty items, 92 and 65 of 140, and together 112, at the auditor's 2.9\% flag rate on clean data
(post hoc, counting only findings about the injected fault: 61 and 110). On scientific code with no
tests shown, reframing the task around the step rather than the whole problem cut the auditor's flags
on correct code from 68.7\% to 36.0\% and its recall from 87.7\% to 72.8\%, on the same programs.

\textbf{Contributions.}
\begin{itemize}\itemsep1pt
\item A measurement programme for an automated auditor's limit against executed ground truth, with
  union recall always reported beside the false-positive rate it cost: controlled changes to the
  number of readings and the rulebook, route comparisons that change model and sampling together,
  re-grading under decision rules, and a post hoc association with what the specification determines.
\item Evidence that no single factor fixes the limit on general-purpose code: each of these moves it,
  and a recall figure without its operating point does not transfer.
\item Operating points of an LLM auditor on scientific code, and comparisons of the LLM tier with
  deterministic verification on reported results (the shipped provenance checks and re-execution) and
  data (a documentation-derived validator), against executed or constructed ground truth. The tiers
  flag different faults, and each fails where the other does not.
\item A methodological finding: in all three science studies the auditor flagged defects in our own
  construction of items meant to be clean, so an evaluation of an auditor must establish that its
  clean items are clean by the auditor's standard, or its false-positive rate measures the evaluator.
\end{itemize}

\textbf{Preregistration and review.} Every study that made model calls was preregistered before its
first one, with hypotheses, intervals and stopping rules fixed in advance, and planned departures
recorded as numbered amendments. Two analyses were not registered in advance, and the paper names
both (Appendix~\ref{sec:method}). Four preregistered hypotheses were killed by their own rules, and
in the science studies one more. Every report was reviewed by a model from a different vendor with
access to the raw records, at between three and twenty-one rounds each, and the reviews are archived
verbatim. Section~\ref{sec:discussion} discusses what the measurements mean. The full method, results
and discussion are in Appendices~\ref{sec:method}--\ref{sec:discussion-full}, and the records, reviews
and the command that regenerates every figure are described in Appendix~\ref{sec:repro}.
